% 分类目录：ml
\documentclass[12pt, a4paper]{article}
\usepackage[margin=2.5cm]{geometry}
\usepackage{ctex}
\usepackage{amsmath, amssymb}
\usepackage{listings}
\usepackage{xcolor}
\usepackage{tabularx}
\usepackage{hyperref}

\lstset{
    language=Python,
    basicstyle=\ttfamily\small,
    keywordstyle=\color{blue},
    commentstyle=\color{green!50!black},
    stringstyle=\color{red},
    showstringspaces=false,
    breaklines=true,
    numbers=left,
    numberstyle=\tiny\color{gray},
    frame=single
}

\title{知识点：三大抽样分布 ($\chi^2, t, F$)}
\author{}
\date{}

\begin{document}
\maketitle

\section{核心定义}
\textbf{中文定义}：抽样分布（Sampling Distribution）是样本统计量的概率分布。在正态总体假设下，衍生出了统计学中最重要的三大分布：卡方分布 ($\chi^2$)、学生 $t$ 分布和 $F$ 分布。它们是假设检验和区间估计的数学基石。

\textbf{英文定义}：Sampling distributions are probability distributions of sample statistics. Under the assumption of a normal population, three fundamental distributions arise: the Chi-square ($\chi^2$) distribution, the Student's $t$ distribution, and the $F$ distribution. They form the mathematical foundation for hypothesis testing and confidence intervals.

\section{卡方分布 ($\chi^2$ Distribution)}

\subsection{定义与构造}
设 $X_1, X_2, \dots, X_n$ 是来自标准正态分布 $N(0, 1)$ 的独立随机变量，则称随机变量
$$ \chi^2 = X_1^2 + X_2^2 + \dots + X_n^2 $$
服从自由度为 $n$ 的卡方分布，记作 $\chi^2 \sim \chi^2(n)$。

\subsection{重要性质与证明}
\begin{itemize}
    \item \textbf{均值与方差}：$E[\chi^2] = n$，$\text{Var}(\chi^2) = 2n$。
    \item \textbf{证明思路}：利用 $X \sim N(0, 1)$ 则 $E[X^2] = 1$，$E[X^4] = 3$。那么 $E[\sum X_i^2] = \sum E[X_i^2] = n$。$\text{Var}(\sum X_i^2) = \sum \text{Var}(X_i^2) = \sum (E[X_i^4] - (E[X_i^2])^2) = \sum (3-1) = 2n$。
    \item \textbf{可加性}：若 $U \sim \chi^2(n_1)$，$V \sim \chi^2(n_2)$ 且相互独立，则 $U+V \sim \chi^2(n_1+n_2)$。
\end{itemize}

\subsection{正态总体样本方差的分布}
\textbf{定理}：设样本来自 $N(\mu, \sigma^2)$，则样本均值 $\bar{X}$ 与样本方差 $S^2$ 相互独立，且：
$$ \frac{(n-1)S^2}{\sigma^2} \sim \chi^2(n-1) $$
\textbf{证明要点}：利用正交变换或分解定理。设 $Z_i = (X_i - \mu)/\sigma$，则 $\sum Z_i^2 = \sum (Z_i - \bar{Z})^2 + n\bar{Z}^2$。其中 $\sum (Z_i - \bar{Z})^2 = (n-1)S^2/\sigma^2$，且 $n\bar{Z}^2 \sim \chi^2(1)$。由 Cochran 定理可知两者独立，故自由度为 $n-1$。

\section{学生 $t$ 分布 (Student's $t$)}

\subsection{定义}
设 $Z \sim N(0, 1)$，$Y \sim \chi^2(n)$，且 $Z$ 与 $Y$ 相互独立，则称随机变量
$$ t = \frac{Z}{\sqrt{Y/n}} $$
服从自由度为 $n$ 的 $t$ 分布，记作 $t \sim t(n)$。

\subsection{性质与应用}
\begin{enumerate}
    \item \textbf{对称性}：关于 0 对称。当 $n \to \infty$ 时，$t$ 分布趋于标准正态分布。
    \item \textbf{证明其在均值检验中的作用}：
    对于 $N(\mu, \sigma^2)$，$\frac{\bar{X} - \mu}{\sigma/\sqrt{n}} \sim N(0, 1)$。由于 $\sigma$ 未知，用 $S$ 代替：
    $$ \frac{\bar{X} - \mu}{S/\sqrt{n}} = \frac{(\bar{X} - \mu)/(\sigma/\sqrt{n})}{\sqrt{S^2/\sigma^2}} = \frac{Z}{\sqrt{\chi^2(n-1)/(n-1)}} \sim t(n-1) $$
\end{enumerate}

\section{$F$ 分布 (F-Distribution)}

\subsection{定义}
设 $U \sim \chi^2(n_1)$，$V \sim \chi^2(n_2)$，且 $U$ 与 $V$ 相互独立，则随机变量
$$ F = \frac{U/n_1}{V/n_2} $$
服从自由度为 $(n_1, n_2)$ 的 $F$ 分布，记作 $F \sim F(n_1, n_2)$。

\subsection{重要性质}
\begin{itemize}
    \item \textbf{倒数性质}：若 $F \sim F(n_1, n_2)$，则 $1/F \sim F(n_2, n_1)$。
    \item \textbf{与 $t$ 分布的关系}：若 $t \sim t(n)$，则 $t^2 \sim F(1, n)$。
    \item \textbf{应用}：方差分析 (ANOVA) 和回归方程的显著性检验。
\end{itemize}

\section{三大分布关系总结}
\begin{table}[h]
\centering
\begin{tabularx}{\textwidth}{|l|X|X|}
\hline
\textbf{分布} & \textbf{构造基础} & \textbf{典型应用} \\
\hline
$\chi^2$ & 标准正态的平方和 & 拟合优度、独立性检验、方差区间估计 \\
\hline
$t$ & 正态 / $\sqrt{\chi^2/n}$ & 均值检验 ($\sigma$ 未知)、回归系数检验 \\
\hline
$F$ & $(\chi^2/n_1) / (\chi^2/n_2)$ & 比较两组方差、方差分析、回归方程检验 \\
\hline
\end{tabularx}
\end{table}

\section{关键代码片段 (Python)}
\begin{lstlisting}
import numpy as np
from scipy import stats

# 1. 卡方分布临界值
alpha = 0.05
df = 10
critical_chi2 = stats.chi2.ppf(1 - alpha, df)

# 2. t 分布概率计算
p_val = stats.t.sf(2.26, df=9) # Survival function (1-CDF)

# 3. F 分布临界值
f_stat = stats.f.ppf(0.95, dfn=5, dfd=10)

print(f"Chi2 Critical: {critical_chi2}")
print(f"P-value (t): {p_val}")
\end{lstlisting}

\section{深度面试题}

\textbf{问：为什么在大样本下不再强求 $t$ 检验而可以使用 $Z$ 检验？}\\
答：根据中心极限定理，样本均值趋向正态分布。同时，根据一致性，当 $n$ 很大时，样本标准差 $s$ 极度接近总体标准差 $\sigma$。此时 $t$ 分布的厚尾特性消失，数值上极度接近标准正态分布，故可用 $Z$ 替代。

\textbf{问：如何证明 $\bar{X}$ 与 $S^2$ 相互独立？}\\
答：这是一个深刻的结论，通常在正态总体假设下成立。数学证明可以通过构造一个正态向量 $\mathbf{X}$，并应用正交变换 $\mathbf{Y} = \mathbf{A}\mathbf{X}$。选取正交矩阵 $\mathbf{A}$ 使得其第一行对应于均值的方向，则变换后的变量 $Y_1$ 正比于 $\bar{X}$，而其余 $Y_2, \dots, Y_n$ 决定了 $S^2$。由于正交变换保持独立性（在正态分布下），故 $\bar{X}$ 与 $S^2$ 独立。

\section{参考文献}
\begin{enumerate}
    \item Cochran, W. G. (1934). \textit{The Distribution of Quadratic Forms in a Normal System, with Applications to the Analysis of Variance}.
    \item Fisher, R. A. (1922). \textit{On the Mathematical Foundations of Theoretical Statistics}.
    \item Mood, A. M., et al. (1974). \textit{Introduction to the Theory of Statistics}. McGraw-Hill.
\end{enumerate}

\end{document}
